\documentclass[11pt]{article}

\usepackage[margin=1in]{geometry}
\usepackage{graphicx}
\usepackage{amsmath, amssymb}
\usepackage{booktabs}
\usepackage{hyperref}
\usepackage{float}
\usepackage{enumitem}
\usepackage{caption}

\hypersetup{colorlinks=true, linkcolor=blue, citecolor=blue, urlcolor=blue}

\title{ECE5242 Project 2: Gesture Recognition \\with Hidden Markov Models}
\author{Dean Xu (hx332)}
\date{March 9, 2026}

\begin{document}
\maketitle

%------------------------------------------------------------------
\section{Introduction}

The goal of this project is to recognize 6 types of hand gestures---Wave, Infinity, Eight, Circle, Beat3, Beat4---from IMU sensor data (3-axis gyroscope + 3-axis accelerometer, so 6D per timestep).

The basic idea is to train one HMM per gesture, and at test time, feed the unknown sequence into all 6 models, compute the log-likelihood for each, and pick the highest one. The whole pipeline has four parts: data preprocessing and vector quantization, HMM algorithm implementation, training and cross-validation, and inference.

Final result on Leave-One-Out cross-validation: \textbf{100\% Top-1 accuracy} (36/36) after adding a feature normalization step that removes device orientation bias. On the 8-file blind test set, predictions show good confidence with large LL margins for most gestures; the main ambiguity is between beat3 and beat4, which makes sense since they look really similar.


%------------------------------------------------------------------
\section{Understanding the Data and the First Big Decision}

\subsection{Data Layout}

There are two sets of training data:

\begin{itemize}[nosep]
    \item \textbf{Repeated gesture}: each file has the same gesture repeated 3--6 times back to back, so the sequences are pretty long (2000--3400 timesteps).
    \item \textbf{Single gesture}: one gesture per file, around 400--850 timesteps, same format as the test data.
\end{itemize}

Each gesture type has 5 repeated files and 1 single file. Each timestep is a 7-dim vector (timestamp + 6D IMU). I drop the timestamp column and just keep the 6 sensor channels.

\subsection{Looking at the Signals}

First thing I did was plot the raw IMU signals for all 6 gestures (Figure~\ref{fig:raw_signals}). The patterns are pretty different visually---wave has this back-and-forth oscillation, circle is smooth and periodic, beat3/beat4 have sharp pulse-like patterns. So yeah, the 6D data clearly carries enough info to tell them apart. HMM should work.

\begin{figure}[H]
    \centering
    \includegraphics[width=0.95\textwidth]{raw_signal.png}
    \caption{Raw IMU signals for the 6 gesture types (one repeated file each, all 6 channels overlaid).}
    \label{fig:raw_signals}
\end{figure}

\subsection{Trying to Segment---and Failing}

Since the repeated files have multiple gestures in them, the obvious idea is to chop them up into individual ones before training. I tried energy-based segmentation:

\begin{enumerate}[nosep]
    \item \textbf{Full 6D energy} (L2 norm of all channels): The baseline energy was always above 10 because the accelerometer picks up gravity ($\sim$9.8~m/s$^2$), so it never drops near zero even when the arm is still.
    \item \textbf{Gyroscope-only energy} (L2 norm of first 3 channels): In theory angular velocity should be zero when still. But in practice the person barely pauses between repetitions---the gyro energy hovers around 0.8--2.5 the whole time, no reliable ``quiet zone'' to cut at.
\end{enumerate}

The results were terrible: 5 files gave me only 5--6 big chunks (basically one chunk per file), when I expected 20+ individual gestures. So segmentation didn't work.

\subsection{The Decision: Cyclic Left-to-Right HMM}

Since I couldn't segment, I changed approach: just use the entire repeated file as one training sequence, and use a \textbf{cyclic left-to-right HMM}.

A standard left-to-right topology only lets states go forward (stay or advance one step), which models a gesture going through a series of phases from start to finish. Adding the cyclic part means the last state can jump back to the first state, representing ``one repetition done, starting the next one''. This way even if there's 3--6 repeats in the sequence, the HMM can handle it naturally.

This decision came straight from Tip~10 in the project PDF: if you don't split the repeated gestures, you can use a left-to-right structure that allows transition from the last state back to the first.

\textbf{Data split}: All 6 files per gesture (5 repeated + 1 single) go into training. The single file is the same format as test data, so including it gives the HMM some exposure to shorter single-gesture sequences. Generalization is checked through Leave-One-Out cross-validation on all 6 sequences (Section~\ref{sec:loocv}), which properly excludes the held-out sample from training each fold.


%------------------------------------------------------------------
\section{Vector Quantization}

\subsection{Why Quantize}

I'm implementing a discrete HMM, so the emission matrix $B$ is an $N \times M$ probability table ($B[j,k] = P(o_t = k \mid s_t = j)$). The observations need to be integers from a finite set. So I need to map the continuous 6D IMU vectors to discrete labels first.

\subsection{Feature Normalization}

Before running K-means, I noticed a big problem with the raw features: the accelerometer channels have way higher variance than the gyroscope channels because of the gravity component ($\sim$9.8~m/s$^2$). Looking at the raw data, about 85\% of the total variance comes from accelerometer and only 15\% from gyroscope. If you run K-means directly on this, the clusters are basically determined by accelerometer values alone---gyroscope information gets ignored.

Even worse, the gravity vector projects differently onto the sensor axes depending on how you hold the phone. So the same gesture performed with a slightly different device orientation gives very different raw accel values, which means the K-means quantization becomes orientation-dependent.

The fix is two-step normalization:
\begin{enumerate}[nosep]
    \item \textbf{Per-sequence mean removal}: Subtract the mean of each sequence. This removes the static gravity bias (which is roughly constant within one recording).
    \item \textbf{Global std normalization}: Divide by the global standard deviation computed across all mean-centered training data. This equalizes the contribution of gyroscope and accelerometer channels.
\end{enumerate}

After normalization, variance contribution becomes roughly 50/50 between gyro and accel. The normalization parameters (global std) are saved and applied identically at test time.

\subsection{K-means Clustering}

After normalization, I pool all 6D vectors from the 36 training sequences and run K-means with $M=70$ clusters, then replace each timestep's normalized 6D vector with its nearest cluster index.

\subsection{Why $M=70$}

The course recommends $M$ between 50--100. I started with 70, checked the clustering quality: no empty clusters, smallest cluster has 272 samples, largest has 4899, distribution looks reasonable (Figure~\ref{fig:cluster_hist}). Later grid search (see Section~\ref{sec:hyperparam}) over $M \in \{65, 70, 75\}$ confirmed 70 works best. Too small loses information, too large means each cluster has too few samples and $B$ estimation gets noisy.

\begin{figure}[H]
    \centering
    \includegraphics[width=0.7\textwidth]{cluster.png}
    \caption{Cluster size distribution for K-means with $M=70$.}
    \label{fig:cluster_hist}
\end{figure}

After quantization each sequence becomes an integer array, same length as before. Each gesture has 6 training sequences (5 repeated + 1 single).


%------------------------------------------------------------------
\section{HMM Implementation and Correctness Checks}

\subsection{The Core}

The HMM class is about 110 lines, with four main methods: \texttt{forward()}, \texttt{backward()}, \texttt{log\_likelihood()}, and \texttt{train()}. It supports both ergodic and left-to-right-cyclic topologies.

\subsection{Numerical Stability: Rabiner Scaling}

This was probably the most important technical decision. In the forward algorithm, $\alpha[i,t]$ is a product of probabilities. With sequences 2000--3000 steps long, $\alpha$ values quickly underflow to 0 (double precision bottoms out around $10^{-308}$), and then everything becomes NaN.

I went with the scaling procedure from Section~V of the Rabiner paper~\cite{rabiner1989}: at each step, normalize $\alpha$ so it sums to 1, and save the scaling factor $c[t]$. The final log-likelihood is:
\begin{equation}
    \log P(O \mid \lambda) = \sum_{t} \log c[t]
\end{equation}
The scaling factors themselves encode all the probability info, no need to sum $\alpha$ at the end. The backward pass uses the same $c[t]$ values so forward and backward stay numerically consistent.

Why scaling instead of log-space: the scaling approach keeps the formulas looking almost the same as the original forward/backward equations, so it's easier to implement. Log-space would mean turning every multiplication into addition and every summation into log-sum-exp, and computing $\xi$ gets messier.

\subsection{Dealing with Zero Emission Probabilities}

If some state never emits a particular observation during training, then $B$ has a zero in that spot. At test time, if that observation shows up, the whole sequence probability goes to zero. And you can't compare two models that both give zero.

Fix: additive smoothing on $B$---add $\varepsilon = 10^{-8}$ to every entry, then renormalize. Intuitively this means the model says ``this observation is extremely unlikely'' instead of ``this observation is impossible''. This needs to happen both at initialization and after every M-step update.

\subsection{Parameter Initialization}

\textbf{$A$ matrix} (transition probabilities): For the left-to-right cyclic topology, only the diagonal (self-loop), super-diagonal (advance), and the cyclic position (last $\to$ first) are nonzero. Self-loop probabilities are initialized to random values between 0.5--0.7 (so each state stays for a few steps on average, covering a chunk of the observation), and the advance probability is whatever's left. An $A_{\text{mask}}$ matrix forces all disallowed positions to zero after every M-step, so the LR topology can't get corrupted during training. Also worth noting: positions that start at zero in $A$ will stay zero forever through EM---this is mathematically guaranteed by the left-to-right constraint.

\textbf{$B$ matrix} (emission probabilities): Sampled from $\text{Dirichlet}(1,1,\ldots,1)$ (uniform random on the $M$-dim simplex), then add $\varepsilon = 10^{-8}$ and renormalize. I don't use uniform initialization (every row $= 1/M$) because symmetric initial conditions cause all states to learn the same $B$ row, and the model basically can't learn anything useful.

\textbf{$\pi$} (initial state distribution): Fixed at $[1, 0, \ldots, 0]$. In a left-to-right model the sequence always starts at state 0, so $\pi$ doesn't need learning.

\subsection{Speed}

The forward recursion core $\alpha_{t} = (A^\top \alpha_{t-1}) \odot B_{:,o_t}$ is done with numpy matrix multiply in one shot, replacing the naive double for-loop. With $N=15$ the difference isn't extreme, but over 2000+ timesteps it adds up a lot.

For $\xi$ computation, I accumulate step by step: compute the $N \times N$ matrix at each timestep and add it directly to \texttt{xi\_sum}, instead of storing the full $(N, N, T-1)$ 3D array. For a sequence with $T=3000$ and $N=15$, this reduces memory from about 5~MB to 1.8~KB. This turned out to be critical when running 30+ consecutive HMM trainings during LOOCV---without it I was getting MemoryError.

\subsection{Toy Data Sanity Check}

Before touching real data, I validated the implementation on a simple problem with a known answer.

\textbf{Setup}: 3000-step sequence---first 1000 steps in state~0 (only emits obs=0), middle 1000 steps in state~1 (emits 0 or 1 with equal probability), last 1000 steps back to state~0. Using ergodic topology here (because the states go $0 \to 1 \to 0$, which needs backward transitions).

\textbf{Expected}:
\[
A \approx \begin{bmatrix} 0.999 & 0.001 \\ 0.001 & 0.999 \end{bmatrix}, \quad
B \approx \begin{bmatrix} 1.0 & 0.0 \\ 0.5 & 0.5 \end{bmatrix}
\]

\textbf{Result}: Converged in 19 epochs, LL $= -707.50$. The learned $A$ and $B$ match the theoretical values closely. LL curve is strictly monotonically increasing (Figure~\ref{fig:toy_ll}), confirming the Baum-Welch implementation is correct. The fact that LL never decreases also implicitly verifies forward-backward numerical consistency---if they were inconsistent, $\gamma = \alpha \times \beta$ normalization would go wrong and LL couldn't stay monotonic.

\begin{figure}[H]
    \centering
    \includegraphics[width=0.7\textwidth]{Toy_LL.png}
    \caption{Toy data training LL curve (strictly increasing, converges in 19 epochs).}
    \label{fig:toy_ll}
\end{figure}

\textbf{Gotcha}: With seed=42, EM got stuck in a local optimum where both states learned nearly identical $B$ rows---they couldn't differentiate. Switched to seed=4 and it converged correctly. This confirmed that EM is sensitive to initialization and you really need to try multiple random seeds.


%------------------------------------------------------------------
\section{Training and Validation}

\subsection{Training Setup}

Hyperparameters: $N=15$ hidden states, $M=70$ observation symbols, Left-to-Right Cyclic topology, seed=42, max 200 epochs, convergence threshold $\text{tol} = 10^{-2}$ (stop when absolute LL change between consecutive epochs drops below 0.01).

I train one independent HMM for each of the 6 gestures. Each model uses that gesture's 6 sequences (5 repeated + 1 single) with multi-sequence Baum-Welch: run forward/backward independently on each sequence (E-step), accumulate $\gamma$ and $\xi$ statistics across sequences, then do a single M-step to update parameters.

\subsection{Training Behavior}

All 6 models converged normally with monotonically non-decreasing LL (Figure~\ref{fig:ll_curves}). Convergence speed varied quite a bit:

\begin{table}[H]
\centering
\caption{Training convergence summary for 6 gesture HMMs.}
\label{tab:training}
\begin{tabular}{@{}lcrl@{}}
\toprule
Gesture & Epochs & Final LL & Notes \\
\midrule
wave   & 108 & $-10451$  & Slowest, several small jumps \\
inf    & 107  & $-13628$ & Also slow, many small improvements \\
eight  & 44  & $-13912$ & Fast convergence \\
circle & 83  & $-6787$  & Highest LL, shortest \& most regular \\
beat3  & 84 & $-10073$ & Big LL jump $\sim$epoch 50 ($\approx$+500) \\
beat4  & 105  & $-14102$ & Large jump around epoch 90 \\
\bottomrule
\end{tabular}
\end{table}

The beat3 LL curve had a big jump around epoch 50---basically EM escaped from one local optimum to a better one. beat4 had a similar jump around epoch 90. This isn't unusual. What matters is the jump is always upward (LL increases), which is consistent with the EM theoretical guarantee.

\begin{figure}[H]
    \centering
    \includegraphics[width=0.95\textwidth]{ll_curves_all.png}
    \caption{Training LL curves for all 6 gesture HMMs. All curves are monotonically non-decreasing, confirming correct Baum-Welch implementation.}
    \label{fig:ll_curves}
\end{figure}

\subsection{Avoiding Data Leakage}

A common pitfall is validating on data that was also used for training. Since I'm using all 6 sequences per gesture for training, I need a proper held-out evaluation scheme. LOOCV handles this: each fold strictly excludes one sequence from that gesture's training, so the held-out sample is never seen during parameter estimation.

Note that the K-means step does see all data (including the held-out sequence), but this is just a fixed preprocessing mapping from continuous features to discrete labels---it doesn't affect HMM parameter learning ($A$, $B$, $\pi$), so it's not leakage.

\subsection{Leave-One-Out Cross-Validation}
\label{sec:loocv}

To get a real sense of generalization, I ran LOOCV: for each gesture's 6 training sequences, leave one out, retrain that gesture's HMM on the remaining 5, then score the left-out sequence against all 6 models and predict with the highest LL. The other 5 gesture HMMs don't need retraining (their training data isn't affected). 6 gestures $\times$ 6 folds = 36 evaluation samples total.

\textbf{Results}:
\begin{itemize}[nosep]
    \item \textbf{Top-1: 36/36 = 100\%}
    \item \textbf{Top-3: 36/36 = 100\%}
\end{itemize}

\begin{figure}[H]
    \centering
    \includegraphics[width=0.55\textwidth]{loocv_confusion.png}
    \caption{LOOCV confusion matrix. Perfect diagonal---all 36 samples correctly classified.}
    \label{fig:confusion}
\end{figure}

Before adding the normalization step (Section~3.2), LOOCV was 35/36 = 97.2\%, with the only error being beat3 [rep4] misclassified as beat4. The normalization fixed this by removing device orientation bias that made the two beating gestures' cluster assignments overlap more than they should.


%------------------------------------------------------------------
\section{Hyperparameter Sensitivity}
\label{sec:hyperparam}

To check if $N=15$ and $M=70$ are actually in the sweet spot, I did a small grid search: $M \in \{65, 75\}$, $N \in \{10, 20\}$, using single holdout as a quick eval metric (train on 5 repeated, validate on 1 single). The metric is correct count + average LL margin (difference between the correct gesture's LL and the runner-up).

\begin{table}[H]
\centering
\caption{Hyperparameter grid search results.}
\label{tab:grid}
\begin{tabular}{@{}cccccl@{}}
\toprule
$M$ & $N$ & Correct & Avg Margin & beat3 Margin & Notes \\
\midrule
70 & 15 & 6/6 & $\sim$2500 & $\sim$200 & baseline \\
65 & 10 & 6/6 & 2502 & 239 &  \\
65 & 20 & 6/6 & 2604 & 161 &  \\
75 & 10 & 6/6 & 2822 & 397 & highest beat3 margin \\
75 & 20 & 6/6 & 2846 & 38.5 & beat3 margin drops hard \\
\bottomrule
\end{tabular}
\end{table}

Key findings:
\begin{enumerate}[nosep]
    \item All combos got 6/6---the model isn't super sensitive to small $N$/$M$ changes.
    \item $N=20$ makes beat3 margin drop from $\sim$200 to 38.5---too many hidden states leads to overfitting, the beat3 and beat4 models start looking too similar.
    \item $M=75$ has slightly better overall margin than $M=65$, but not by much.
\end{enumerate}

\textbf{Conclusion}: Stick with $N=15$, $M=70$. Don't go too high on $N$ (overfitting), $M$ is fine anywhere in 65--75.


%------------------------------------------------------------------
\section{Final Model and Test Results}

\subsection{Final Model}

Since training already uses all 6 sequences per gesture, the models from the main training run are the final models. They're saved to the \texttt{models/} directory along with the K-means model and normalization parameters.

\subsection{Inference Pipeline}

The inference script \texttt{test\_classifier.py} is separate from the notebook, about 50 lines:

\begin{enumerate}[nosep]
    \item Load K-means model (70 cluster centers) + 6 HMM models ($A$, $B$, $\pi$ matrices) + normalization parameters (global std)
    \item Read test file $\to$ drop timestamp column $\to$ get 6D IMU data
    \item Normalize: subtract per-sequence mean, divide by global std (same as training)
    \item K-means predict $\to$ quantize the normalized 6D vectors into integer labels
    \item Run scaled forward on all 6 HMMs $\to$ get 6 log-likelihoods
    \item Highest LL = Top-1 prediction; also output Top-3 with their LL values
\end{enumerate}

\subsection{Test Set Results}

The test set has 8 unlabeled gesture files. I ran \texttt{test\_classifier.py} with the final models:

\begin{table}[H]
\centering
\caption{Test set classification results (Top-3 predictions with log-likelihoods).}
\label{tab:test}
\begin{tabular}{@{}lllcrrr@{}}
\toprule
File & Top-1 & Top-2 & Top-3 & LL-1 & LL-2 & LL-3 \\
\midrule
test1.txt & inf    & eight & beat3  & $-749$   & $-6833$  & $-7018$ \\
test2.txt & beat3  & beat4 & inf    & $-797$   & $-1101$  & $-7033$ \\
test3.txt & inf    & eight & beat3  & $-647$   & $-5376$  & $-7492$ \\
test4.txt & beat4  & beat3 & inf    & $-686$   & $-847$   & $-7532$ \\
test5.txt & circle & beat4 & beat3  & $-326$   & $-6216$  & $-6430$ \\
test6.txt & eight  & inf   & wave   & $-1745$  & $-4066$  & $-7487$ \\
test7.txt & wave   & eight & inf    & $-338$   & $-3409$  & $-5036$ \\
test8.txt & beat3  & beat4 & circle & $-940$   & $-1009$  & $-7700$ \\
\bottomrule
\end{tabular}
\end{table}

\textbf{Analysis}:
\begin{itemize}[nosep]
    \item \textbf{Top-1 distribution}: inf $\times$2, beat3 $\times$2, beat4 $\times$1, circle $\times$1, eight $\times$1, wave $\times$1. All 6 gesture types appear in the predictions, which looks reasonable.
    \item \textbf{LL margins}: Most predictions have large Top-1 vs Top-2 gaps (e.g., test5 circle: margin $>$5800, test3 inf: margin $>$4700), suggesting high confidence. test4 and test8 have smaller margins (161 and 69 respectively), both involving beat3/beat4 ambiguity---consistent with LOOCV showing these two are the hardest pair.
    \item \textbf{Normalization effect}: Before normalization, the Top-1 predictions were heavily skewed toward a few gesture types. After normalization, the predictions are spread across all 6 gestures with much better LL separation.
\end{itemize}


%------------------------------------------------------------------
\section{Discussion}

\subsection{Why beat3 and beat4 Get Confused}

Both gestures are basically up-and-down beating motions. The main difference in the IMU signal is just temporal---3 beats vs 4 beats---not spatial. A fixed-state-count LR-HMM has a hard time precisely capturing the difference between ``repeat 3 times'' and ``repeat 4 times'' because that requires very fine-grained temporal modeling.

On the training data (LOOCV), normalization completely fixed the beat3/beat4 confusion. But on test data, the small LL margins between beat3 and beat4 for several files (e.g., test4: margin = 161, test8: margin = 69) suggest these two remain hard to separate. This makes sense---normalization helps a lot but the fundamental similarity remains. If I wanted to improve further: (a)~use more hidden states specifically for beat3/beat4; (b)~add time-domain features like first differences or short-time energy; (c)~try multiple random seeds and keep the best model.

\subsection{Limitations}

Discrete HMM + K-means quantization is what the course recommends, and it's simple and works pretty well, but there are inherent limitations. K-means quantization inevitably loses information---different gestures' IMU vectors might get mapped to the same cluster. A continuous HMM (with GMM emission probabilities) could avoid this quantization loss, but it's way more complex to implement.

\subsection{Engineering Lessons}

\begin{itemize}[nosep]
    \item \textbf{Memory}: Initially Baum-Welch stored the full $\xi$ 3D array $(N, N, T{-}1)$. When running LOOCV with 30+ consecutive trainings this hit MemoryError. Switching to step-by-step accumulation of the $(N, N)$ statistics fixed it.
    \item \textbf{Batch execution}: The long LOOCV training was done through a standalone Python script with checkpoint saving and real-time progress output, instead of running the whole notebook as a black box. Much more controllable.
    \item \textbf{Feature normalization matters}: Without normalization, LOOCV was 97.2\% and test accuracy was poor---the model was overfitting to device orientation. Per-sequence mean removal + global std normalization fixed this and brought LOOCV to 100\%. The lesson: always think about what your features actually encode, not just the model architecture.
\end{itemize}


%------------------------------------------------------------------
\begin{thebibliography}{9}
\bibitem{rabiner1989}
L.~R. Rabiner, ``A Tutorial on Hidden Markov Models and Selected Applications in Speech Recognition,'' \textit{Proceedings of the IEEE}, vol.~77, no.~2, pp.~257--286, 1989.
\end{thebibliography}

\end{document}
